\scanpage{14}
\fbox{\S 1} \uline{Préliminaires}
\margin{$\longleftarrow\quad 1,\unsure{2}\quad !!\longrightarrow$}
% slanted
\margin{\uline{NB.} Les préliminaires se \unsure{compliquent|complètent} \unsure{et se poursuivent} \unsure{dans le paragr. 2 jusqu'en} \illegible{1} \unsure{et|p.} \unsure{$\leftarrow$}}
% slanted
\margin{\deleted{\unsure{Il suffit que} \unsure{il soit défini pour les} limites \unsure{finies}, \unsure{sommes} \unsure{et} \illegible{2} \unsure{deux objets}}}

\uline{Proposition 1} Soit $M$ une catégorie stable \unsure{pour} $\unsure{\varinjlim}$, \unsure{quelconque|quelconques}, $C\subset\mathrm{Fl}\,M$ contenant les iso, stable par \unsure{cochgt.|chgt.} de base et par composition \uline{transfinie}. Soit, pour tout \inserted{petit} ens. ordonné $I$, $C_I\subset\mathrm{Fl}\, M(I^\circ)$ l'ens des flèches \inserted{$u_.:X_.\unsure{-|\to}Y_.$} telles que $\forall i_0\in I$, \deleted{considérant} \inserted{dans} le diagramme

\diagram{\begin{tikzcd}[ampersand replacement=\&,column sep=large,row sep=large]
 \&X'_{i_0}=\varinjlim_{i<i_0}X_i
 \arrow[r]\arrow[d,"{\alpha_{i_0}=\lim_{i<i_0}u_i}"']\arrow[dr,phantom,"\text{cocart}"]
 \& X_{i_0}\arrow[d,"{\unsure{\beta}_{i_0}}"']\arrow[dd,bend left=55,dashed,"{u_{i_0}}"]\\
 (\alpha)\&Y'_{i_0}=\underset{i<i_0}{\underline{\lim}}\,Y_i\arrow[r]
 \&\widetilde{Y}_{i_0}\arrow[d,"{\widetilde{u}_{i_0}}"']\&\text{(défini comme somme amalgamée)}\\
 \&{}\&Y_{i_0}
\end{tikzcd}}
% slanted
\margin{Condition stable par passage de $I$ à un sous-ens. ordonné \uline{ouvert} $I'$ de $I$}

\unsure{la flèche|les flèches} $\widetilde{u}_{i_0}$ \unsure{soit|soient} dans $C$. \deleted{Soit alors $I$ \unsure{soit}} \inserted{Je dis que} \deleted{alors} si $I$ est \uline{bien ordonné}, \deleted{\unsure{et}} alors
$u_.\in C_I$ implique a) que $\varinjlim u_.:\varinjlim X_i-\varinjlim Y_i$ est dans $C$, et b) $\forall i\in\deleted{\unsure{C}} I$, $(u_i:X_i-Y_i)\in C$.

% slanted
\margin{Dire aussi que l'on a b$'$) Si $i'\leq i$, alors \unsure{dans le} diagramme \unsure{similaire|tributaire} à \unsure{$(\alpha)$}
\diagram{\begin{tikzcd}[ampersand replacement=\&,column sep=tiny,row sep=tiny]
 X'_{i'}\arrow[r]\arrow[d] \& X'_i\arrow[d]\\
 Y'_{i'}\arrow[r,-]\arrow[ur,phantom,"\text{cocart}" description] \&\widetilde{Y}'_i\arrow[d,"{u_{i',i}}"]\\
 \&Y'_i
\end{tikzcd}}
on a $u_{i',i}\in C$ \unsure{ce qui} généralise a) \unsure{s'obtient}, retrouve cas en \unsure{prenant} $i'=i_0,\ i=i_0+1)$ le cas b) s'obtient en faisant $i'=0$, $i$ quelconque, et l'hyp. $(X'_{i_0}\to Y'_i)\in C$ en \unsure{supposant} $\widetilde{u}_i:X_i\to Y_i$ \unsure{un} \illegible{1} $i$, \unsure{afin} $\in C$.}

\uline{Cor.} \deleted{\unsure{Démonstration}} Soit $W$ \inserted{$\subset\mathrm{Fl}\,M$, \uline{localiseur}} modérément saturé, tel que $\overset{\inserted{TC\overset{\mathrm{déf}}{=}}}{C\cap W}$ stable par coch. de base, et par composition \uline{transfinie}. Soit $u_.:X_.-Y_.$ dans $C_I$.

Alors $u_.\in W_{I^\circ}$ \deleted{\unsure{(I)}} (i.e. $\forall i\in I$, on a $u_i\in W$) ssi $u_.\in(TC)_I$, i.e. ssi pour \deleted{les} $i\in I$, la flèche $\widetilde{u}_i$ de $(\alpha)$ est dans $TC$, i.e. \deleted{ce qui}
% archive sheet number: 13

